\section{RTT adaptativo --- implementação}

\preencher{Descrever onde e como o estimador é atualizado no código e quais
amostras são ignoradas após retransmissões.}

\subsection{Estimativa EWMA}

Apresentar a fórmula implementada:
\begin{equation}
  \mathrm{EstimatedRTT} \leftarrow
  (1-\alpha)\,\mathrm{EstimatedRTT}+\alpha\,\mathrm{SampleRTT}.
\end{equation}
\preencher{Informar a fórmula completa do timeout, incluindo o desvio do RTT,
e os valores escolhidos para $\alpha$ e $\beta$.}

\subsection{Justificativa dos parâmetros}
\preencher{Justificar $\alpha$ e $\beta$ e mostrar, com dados do experimento,
o impacto dos valores escolhidos.}

\figura{rtt_estimado_vs_medido.png}{0.9\linewidth}
  {RTT estimado versus RTT medido por segmento durante a transferência.}
  {fig:rtt-estimado}

\preencher{Inserir a análise da convergência do estimador, incluindo os trechos
com variação, picos e possíveis perdas.}

\figura{timeout_adaptativo.png}{0.9\linewidth}
  {Evolução do timeout adaptativo ao longo da transferência.}
  {fig:timeout-adaptativo}

\preencher{Explicar como o timeout acompanha as variações de RTT e se houve
estimativas excessivamente conservadoras ou agressivas.}
